\cotefile{46}
\scanpage{24}

\begin{center}
Définition du morphisme Verschiebung.
\end{center}

Je ne sais pas exactement ce que tu veux que je définisse\deleted{\unsure{z|s}}. Cependant, comme tu affirmes avoir perdu mon tapis de l'année dernière, je vais essayer de te le revendre. Pour cela, je vais supposer que le schéma de base est le spectre d'un corps $k$ de caractéristique $p\handwritten{>}0$, parfait ou non : à toi de voir ce qui se passe dans les autres cas.

Je vais commencer par le cas d'un $k$-groupe commutatif \uline{affin\deleted{\unsure{t|é}}e} $G$, d'anneau des coordonnées $\mathbb{A}$. Les notations sont les suivantes :

\noindent - $\mathbb{A}^{(p)}$ est la $k$-algèbre déduite de $\mathbb{A}$ par l'extension des scalaires $F_k:k\mathrel{\text{-{}-{}-{}-{}-}\!\handwritten{>}}k$, où $\mathbb{F}_k(x)=x^p$. Un élément de $\mathbb{A}^{(p)}$ s'écrira donc $x\handwritten{\otimes}a$, où $x\handwritten{\in}k$ et $a\handwritten{\in}\mathbb{A}$ ; on notera la règle de calcul $x\handwritten{\otimes}y.a=x.y^p\handwritten{\otimes}a$. Si $k$ est parfait, $\mathbb{A}^{(p)}$ s'identifie à la $k$-algèbre déduite de $\mathbb{A}$ par la restriction des scalaires $x\mathrel{\text{-{}-{}-{}-{}-}\!\handwritten{>}}x^{1\handwritten{/}p}$.

\noindent - $F:\mathbb{A}^{(p)}\mathrel{\text{-{}-{}-{}-{}-}\!\handwritten{>}}\mathbb{A}$ est l'application \deleted{\illegible{1}} $x\handwritten{\otimes}a\mathrel{\text{-{}-{}-{}-{}-}\!\handwritten{>}}x.a^p$.

\noindent - $G^{(p)}$ est le groupe algébrique d'anneau des coordonnées $A^{(p)}$ ; on désignera encore par $F$ le morphisme de $G$ dans $G^{(p)}$ qui est défini par l'application $F:\mathbb{A}^{(p)}\mathrel{\text{-{}-{}-{}-{}-}\!\handwritten{>}}A$.

\noindent - $\handwritten{\otimes}^{p}A$ désigne le produit tensoriel de $p$ exemplaires de $A$ ; $\handwritten{\Delta}^{p}$ est l'application de $\mathbb{A}$ dans $\handwritten{\otimes}^{p}\mathbb{A}$ qui est définie par la diagonale $\handwritten{\Delta}$ dela coalgèbre $\mathbb{A}$ ; $m^p$ est
% archive sheet number: 23
\scanpage{25}
\mspage{-2}
l'application de $\handwritten{\otimes}^{p}\mathbb{A}$ dans $\mathbb{A}$ qui est définie par la multiplication $m$ de $A$.

\noindent - $\handwritten{\Sigma}^{p}\mathbb{A}$ est la sous-algèbre de $\handwritten{\otimes}^{p}\mathbb{A}$ qui est formée des tenseurs symétriques d'ordre $p$. Comme la diagonale $\handwritten{\Delta}$ est commutative, $\handwritten{\Delta}^{p}$ est le composé d'une application $a:\mathbb{A}\mathrel{\text{-{}-{}-{}-{}-}\!\handwritten{>}}\handwritten{\Sigma}^{p}\mathbb{A}$ et de l'injection canonique $i:\handwritten{\Sigma}^{p}\mathbb{A}\mathrel{\text{-{}-{}-{}-{}-}\!\handwritten{>}}\handwritten{\otimes}^{p}\mathbb{A}$.

Si $(a_i)$ est une base de $\mathbb{A}$, les produits $b=a_{i_1}\handwritten{\otimes}\ldots\handwritten{\otimes}a_{i_r}$ forment une base de $\handwritten{\otimes}^{p}\mathbb{A}$. Si $t_b$ est le sous-groupe du groupe symétrique $s_p$ formé des permutations laissant fixe $b$, nous posons $S(b)=\handwritten{\sum}\handwritten{\sigma}_{\deleted{\unsure{i}}j}(b)$, les $\handwritten{\sigma}_{\deleted{\unsure{i}}j}$ parcourant un système de représentants de $s_p\handwritten{/}t_b$. Il est clair que les $S(b)$ forment une base de $\handwritten{\Sigma}^{p}\mathbb{A}$.

\noindent - $S^{p}\mathbb{A}$ est l'image de $\handwritten{\otimes}^{p}\mathbb{A}$ dans l'algèbre symétrique de $A$. Comme la multiplicatipn $m$ est comm\deleted{\unsure{j|y}}utative, $m^p$ est le composé de la surjection canonique $q:\handwritten{\otimes}^{p}\mathbb{A}\mathrel{\text{-{}-{}-{}-{}-}\!\handwritten{>}}S^p\mathbb{A}$ et d'une appli\deleted{\unsure{s|x}}cation $b:S^p\mathbb{A}\mathrel{\text{-{}-{}-{}-{}-{}-}\!\handwritten{>}}\mathbb{A}$.

Si $(a_i)$ est une base de $\mathbb{A}$, $S^p\mathbb{A}$ est l'ensemble des polynô mes homogènes de degré $p$ en les $a_i$.

\noindent - Nous allons définir une application $r:\handwritten{\Sigma}^{p}\mathbb{A}\mathrel{\text{-{}-{}-{}-{}-}\!\handwritten{>}}\mathbb{A}^{(p)}$ en exhibant les $r(S(b))$ :

si $a_{i_1}=\ldots=a_{i_r}=a$, nous posons $r(S(b))=r(a\handwritten{\otimes}\ldots\handwritten{\otimes}a)=1\handwritten{\otimes}a$.

si les $a_{i_j}$ ne sont pas tous égaux, nous posons $r(S(b))=r(\handwritten{\sum}\handwritten{\sigma}_j(b))=0$.

\noindent - Nous définissons de même une application $j:\mathbb{A}^{(p)}\mathrel{\text{-{}-{}-{}-{}-}\!\handwritten{>}}S^p\mathbb{A}$ telle que
\[
j(1\handwritten{\otimes}a_i)=a_i^p=q(a\handwritten{\otimes}\ldots\handwritten{\otimes}a)\ .
\]
% archive sheet number: 24
\scanpage{26}
\mspage{-\handwritten{3}}

Dans ces conditions on voit facilement que $\unsure{q}.i=j.r$ (voir par exemple Serre, groupes algébriques et corps de classes, preuve du lemme 11, p.\handwritten{62} ; d'est là qu'intervient l'hypothèse que le schéma de base est le spectre d'un corps ; si l'on veut généraliser, on aura donc à étudier \deleted{\unsure{h|m}}le morphisme $q.i$).

On aboutit finalement au diagramme commutatif
\diagram{
\begin{tikzcd}[ampersand replacement=\&,column sep=1.0em,row sep=2.5em]
\mathbb{A}
  \handwritten{\arrow[rr]}
  \arrow[rr,phantom,"\handwritten{\Delta}^{p}"]
  \handwritten{\arrow[dr,"a"]}
\& \& \handwritten{\otimes}^{p}\mathbb{A}
  \handwritten{\arrow[rr]}
  \arrow[rr,phantom,"m^p"]
  \handwritten{\arrow[dr,"q"']}
\& \& \mathbb{A}
\\
\& \handwritten{\Sigma}^{p}\mathbb{A}
  \handwritten{\arrow[ur,"i"']}
  \handwritten{\arrow[dr,"r"]}
\& \& S^p\mathbb{A}
  \handwritten{\arrow[ur,"b"]}
\&
\\
\& \& \mathbb{A}^{(p)}
  \handwritten{\arrow[ur,"j"]}
\& \&
\end{tikzcd}
}

Or $m^p.\handwritten{\Delta}^{p}$ n'est autre que l'homomorphisme d'anneau correspondant à la multiplication par $p$ dans $G$. Cet homomorphisme sera noté $\uline{p}$. De plus, $b.j$ n'est autre que le Frobenius $F$. Si nous notons $V$ et appelons Verschiebung l'application $r.a$, on a $F.V=\uline{p}$. Enfin, comme le diagramme est autodual, $F$ et $V$ s'échangent dans la dualité de Cartier.

Dans le diagramme ci-dessus, $S^p\mathbb{A}$ est le seul espace vectoriel à ne pas être muni d'une structure d'algèbre (c'est une coalgèbre, tandis que $\handwritten{\Sigma}^{p}A$ est une algèbre mais non une coalgèbre). Si \unsure{nons|noms|nous} ne nous intéressons pas à la dualité et si nous voulons raisonner uniquement en termes de variétés \unsure{algébriques|álgébriques}, on obtient donc le diagramme
% archive sheet number: 25
\scanpage{35}

et $H^0$ sont exacts à gauche, on en conclut un homomorphisme de diagrammes exacts d'ensembles
% No arrow strokes are visible between these six objects.
\diagram{
\begin{tikzcd}[ampersand replacement=\&,column sep=2.0em,row sep=1.7em]
H^0(X,F) \& {\unsure{\mathbb{H}|H}}^{0}(X,G) \& H^0(X,H) \\
H^0(Y,F_Y) \& H^0(Y,G_Y) \& H^0(Y,H_Y),
\end{tikzcd}
}
où d'après ce qui précède les flèches verticales sont injectives\inserted{(ennparticulier la dernière}, et où par hypothèse la flèche verticale médiane est bijective. (NB les indices $Y$ désignantt les restrictions à $Y$)? On en conclut formellement que la première flèche verticale est également bijective.
% The inserted parenthesis runs into the right edge; no closing parenthesis is visible.

Nous appliquerons cette remarque au cas \deleted{où}
\[
G=h\ h\ (F)\ x\ (\quad(F))\ ,
\]
% Spaces in this and the following formulas reproduce unfilled places in the typescript.
et à l'homomorphisme évident $F\qquad G$, déduit des homomorphismes naturels dans les deux facteurs. Cet homomorhisme est injectif, comme il résulte immédiatement (par considération des fibres) de l'hypothèse $X=Y\quad(\overline{X})$. Il reste à prouver l'assertion de bijectivité pour $G$, ou ce qui revient au même, pour les deux facteurs, donc pour des faisceaux de la forme $h(\mathbb{H})$ (resp. $(\overline{F})$),où $H$ (resp. $\overline{F}$) est un faisceau sur $Y$ (resp. sur $\overline{X}$). or l'assertion est triviale pour un faisceau de la forme $h(H)$. Pour le cas d'un faisceau de la forme $(\overline{F})$, compte tenu du fait que l'homomorphisme de c\deleted{\unsure{y}}hangement de base $\unsure{H|h}(\quad(\overline{F}))\qquad{}_{Y}(\overline{h}(\overline{F}))$ est bijectif, on voit que l'homomorphisme $H^0(X,\quad(\overline{F}))\qquad H^0(Y,h\quad(\overline{F}))$ s'identifie à l'homomorphisme $H^0(\overline{X},\overline{F})\qquad H^0(\overline{Y},\overline{h}\ (\overline{F}))$, qui est bijectif par hypothèse. Cela achève la démonstration de (i).

(ii) La conclu\deleted{\unsure{t|b}}sion d\deleted{'}e\deleted{\unsure{inject}} bijectivité pour $H^0(X,F)\qquad\mathbb{H}^0(Y,h\ (F))$ \margin{\handwritten{$\unsure{\checkmark|\backslash}$}}pour tout faisceau d'ensembles $F$ sur $X$, déjà obtenue dans (i), assure \deleted{\unsure{par un argument immédiat}} que pour tout faisceau en groupes $F$ sur $X$ (pas nécessairement ind-finis), l\deleted{'}e\deleted{homomorphisme du} foncteur restriction
% archive sheet number: 34
\scanpage{37}

phismes de restriction
\[
H^n(\overline{X},\overline{F})\qquad H^n(\overline{Y},\overline{\unsure{H|h|\mathbb{H}|B}}\ (\overline{F}))\qquad\text{pour }n\quad0\ ,
\]
et les homomorphismes de changement de base
\[
h\ (R^n\ (\overline{F}))\qquad R^n\quad{}_{Y}\ (\overline{h}\ (\overline{F}))\qquad\text{pour }n\quad0,
\]
soient bijectifs. Alors pour tout faisceau abélien de torsion $F$ sur $X$, les homomorphismes de restriction
\[
H^n(X,F)\qquad H^n(Y,h\ (F))\qquad\text{pour }n\quad0
\]
sont \deleted{\unsure{les|des}} bijectifs.
% The spaces between the formula members, and between n and 0, contain no visible signs.

\uline{Démonstration}.\quad (i)\quad Pour l'assertion\deleted{n} non respée, on utilise le carré commutatif (où on pose $\overline{F}=\quad(F)$):
% No arrow strokes are visible in this square.
\diagram{
\begin{tikzcd}[ampersand replacement=\&,column sep=3em,row sep=3em]
H^0(X,F) \& {\unsure{\mathbb{H}|H}}^0(Y,F) \\
H^0(\overline{X},\overline{F}) \& H^0(\overline{Y},\overline{F})\quad,
\end{tikzcd}
}
où par hypothèse la deuxième flèche horizontale est injective, ce qui montre que deux éléments de $H^0(X\deleted{\unsure{n}},F)$ ayant même image dans $H^0(Y,F)$ ont aussi même image dans $H^0(\overline{X},\overline{F})$, \deleted{\unsure{donc|ains}}\inserted{i.e.} ont mêmes images dans les fibres \deleted{géo} de $F$ en les points géométriques au dessus des points de $(\overline{X})$, donc dans \uline{toutes} les fibres puisque $X=Y\quad(\overline{X})$, ce qui prouve qu'elles sont égales. Pour prouver l'assertion respée, soit $F$ un faisceau d'ensembles sur $X$ et prouvons que $H^0(X,F)\qquad H^0(Y,h\ (F))$ est surjectif (donc bijecttif d'après ce qui précède). Montrons qu'il suffit pour ceci de trouver une injection $F\qquad G$, telle que l'assertion analogue soit vraie pour $G$. En effet, d'après les propriétés d'exactitude générales d'une catégorie de faisceaux, $F$ peut alors s'écrire comme un conoyau d'un couple de flèches $G\qquad H$ (où on peut prendre pour $H$ la somme amalgamée $\unsure{\deleted{F}G|\deleted{G}F}\quad{}_{F}G$, si on veut). Tenant compte que le foncteur $H$ est exact\deleted{n}, \deleted{\unsure{on obtient via}} et
% archive sheet number: 36
